% ID: SEL-D0071-OB-E01
% ESTADO: validada. S071.
\textbf{a)} Sean $x,y,z$ los números de latas de A, B y C. El número total, el peso en kg y el coste en euros dan
\[
\begin{cases}
x+y+z=20,\\
0{,}25x+0{,}50y+z=10,\\
x+1{,}8y+3{,}3z=35{,}6.
\end{cases}
\]
\textbf{b)} La segunda ecuación multiplicada por cuatro es $x+2y+4z=40$. Al restar el doble de la primera queda $x=2z$, luego $y=20-3z$. Sustituimos en la ecuación del coste:
\[
2z+1{,}8(20-3z)+3{,}3z=36-0{,}1z=35{,}6.
\]
Sale $z=4$, de donde $x=8$ e $y=8$. Se compraron $\boxed{8,\ 8,\ 4}$ latas de A, B y C, respectivamente. Estos recuentos dan $20$ latas, $10$ kg y $35{,}60$ euros.
